\section{Study Significance}

This thesis intends to contribute to science on multiple frontiers. Methodically, the use of online experiments generally helps to expand the circle of participants and to be able to make more representative statements about the human psyche. 
The preponderance of psychological research continues to draw from a limited and unrepresentative sample—predominantly from WEIRD societies, which has been shown to inadequately reflect the diversity of the global human population.
This is even leveraged by the use of Lucid as an recruitment platform offering a wider range of accessible countries and languages than traditional recruiters. Since this is the first research program to use this recruiter, it is not only taking advantage of the new potentials but also explores the mechanisms of recruitment on this platform. By doing so we set the way for further research exploiting these massive potentials.
In this study, we tested 4,351 participants of 25 languages across three continents Europe, Asia and Africa.

In the domain of color terms, this study closes an important knowledge gap that has existed for a long time. 
Under the assumption that globalization effects diminish differences across languages a lot of research was done investigating unwritten languages of small scale societies.
Conclusions derived from these data sets could only examine the developmental mechanisms in these environments.
Further studies investigated single industrialized languages, but varied in methodology and the availability of data sets and therefore do not allow comparison.
The data set obtained covers a range of 25 mostly industrialized languages and recruited participants fully exposed to cultural exchange through the Internet.
Using the data set obtained, we are now in a position to examine the effects of globalization on the contingent of color terms in industrialized languages.

With growing capabilities, LLMs are becoming more and more important in the daily lives of Internet users around the world. This also increases the importance of understanding how they work. Do they simulate human behavior or operate more independently?
Conducting experiments analogously allows to determine the relationship between human and GPT4 behaviour. Even more importantly, LLMs can help determine whether the perceptual world can be inferred from language.
In the context of this study, GPT4 is considered to be a multilingual agent highly exposed to cultural exchange, functioning as a control group.
